\section{模型假设}

本题给出的计算图、硬件配置与评估程序构成模型的既定条件。计算图中的张量大小、核内计算周期、依赖方向和输入输出位置固定；方案只决定非边界搬运算子的子图归属、子图的核心分配以及每核的子图顺序。原始边界搬运、跨任务交接、容量不足时的换出与取回、流水线占用和共享带宽竞争均按题目规定生成与评价。下列条件分别明确题设执行规则和理论推导的适用范围。

\begin{enumerate}
  \item \textbf{题设执行条件。} 给定同构核心数、私有L1和UB容量及DDR总带宽；场景A与场景B的执行任务组成和等待规则分别固定，第三问另有独立带宽的只读L2。对同一问题、图与配置，所有候选均交由同一官方评估语义判定可行性和完成周期。该条件使不同方案在相同的换出策略、缓存替换和流水线服务规则下比较。
  \item \textbf{固定工作量条件。} 在理论下界的计算部分，将原图中每个待执行计算算子的题给服务周期视为非负且不可省略，并保留其必需的数据依赖。所有合法方案均保留这部分工作量。该条件用于计算路径和每种计算流水线的工作量松弛；实际完成时间仍包含题定搬运、等待和资源争用。
  \item \textbf{必要流量的条件使用。} 只有在某些字节对全部合法方案都必须由DDR服务，且初始缓存状态、逻辑张量别名和输出要求均支持这一判断时，才把这些字节计入全局DDR服务下界。可由只读L2服务的读取按实际服务路径单独分析。必要流量用于下界证明和停止判据，额外COPY则按题定口径统计。
\end{enumerate}

通信量与片上活跃集估计用于候选排序；容量可行性、Spill数量、缓存命中和总体任务执行时间（Makespan）由完整执行评价确定。各问的经验结果均以其对应问题和配置的实际评估响应为准。

